%% =====================================================================
%%  CAPITULO 2 - FUNDAMENTACAO TEORICA
%% =====================================================================

\chapter{Fundamentação Teórica}
\label{cap:fundamentacao}

Este capítulo estabelece os conceitos necessários para compreender o problema e a
solução proposta. A primeira parte trata de engenharia de software --- ciclo de
vida, requisitos, verificação e validação, testes e qualidade. A segunda trata de
agentes de inteligência artificial --- modelos de linguagem, agentes autônomos e
sistemas multiagente. A terceira trata do domínio de aplicação. A seção final
sintetiza as duas áreas e identifica a lacuna que o trabalho aborda.

\section{Engenharia de software e ciclo de vida}
\label{sec:ciclo-de-vida}

Engenharia de software é a aplicação de abordagem sistemática, disciplinada e
quantificável ao desenvolvimento, operação e manutenção de software
(\citeonline{sommerville2019}). A palavra decisiva é \emph{sistemática}: a
diferença entre programar e engenhar software está na existência de um processo
que torne o resultado previsível e inspecionável.

Processos de desenvolvimento organizam-se tradicionalmente em fases ---
elicitação, análise, projeto, implementação, teste, implantação e manutenção ---
ainda que a execução real raramente seja estritamente sequencial
(\citeonline{pressman2016}). Modelos incrementais e iterativos reconhecem que
requisitos mudam e que a retroalimentação é inevitável. O que não se altera entre
modelos é a assimetria de custo: quanto mais tarde um defeito de requisito é
descoberto, maior o custo de corrigi-lo. Essa assimetria pode ser formalizada
como uma progressão geométrica sobre as fases do ciclo de vida:

\begin{equation}
	C(f) = C_0 \cdot \alpha^{\,f-1}, \qquad \alpha > 1
	\label{eq:custo-fase}
\end{equation}

\noindent em que $C(f)$ é o custo de correção quando o defeito é detectado na
fase $f$, $C_0$ é o custo de correção na fase de requisitos e $\alpha$ é o fator
de ampliação por fase. \textbf{O parâmetro $\alpha$ depende do domínio e não foi
medido neste trabalho}; a equação serve para formalizar a direção do efeito, não
para quantificá-lo. A implicação prática é direta: investir em especificação é
investir em reduzir $f$.

\section{Especificação de requisitos}
\label{sec:requisitos}

A norma ISO/IEC/IEEE 29148 define os processos e os produtos da engenharia de
requisitos, incluindo as características que tornam um requisito verificável
(\citeonline{iso29148}). Um requisito bem formado é singular, não ambíguo,
completo, consistente, factível, necessário e --- sobretudo para os fins deste
trabalho --- \emph{verificável}: deve existir um procedimento finito e objetivo
capaz de decidir se o requisito foi atendido.

A literatura separa requisitos em duas classes. Requisitos funcionais descrevem o
que o sistema deve fazer. Requisitos não funcionais descrevem como o sistema deve
se comportar em atributos de qualidade --- desempenho, usabilidade, segurança,
manutenibilidade (\citeonline{sommerville2019}). Requisitos não funcionais são
notoriamente mais difíceis de verificar, porque frequentemente são enunciados
como adjetivos. A prática recomendada é convertê-los em métricas: em vez de
``deve ser rápido'', ``o intervalo entre passos de simulação não deve exceder
$t_{\max}$''.

Uma especificação útil também precisa declarar suas \emph{restrições} ---
decisões já tomadas que o projeto não pode contradizer --- e seus \emph{casos de
borda}, isto é, as entradas ou estados limítrofes que costumam revelar defeitos.
Este trabalho adota a convenção de nomear cada elemento com identificador estável
(\texttt{FR-XXX}, \texttt{NFR-XXX}, \texttt{CA-XXX}), prática que viabiliza a
rastreabilidade automática entre artefatos.

\section{Verificação e validação}
\label{sec:vv}

Verificação e validação são atividades distintas, distinção formalizada por
\citeonline{boehm1984}. A verificação responde à pergunta ``o produto foi
construído corretamente?'', isto é, se ele está em conformidade com a
especificação. A validação responde à pergunta ``foi construído o produto
correto?'', isto é, se ele satisfaz a necessidade que motivou o projeto.

A distinção tem consequência prática imediata para este trabalho. A verificação
de doze critérios de aceitação demonstra conformidade com a especificação; ela
\emph{não} demonstra que o produto é útil, agradável ou que o público-alvo o
adotaria. Afirmações sobre adequação ao uso exigiriam validação com usuários, que
está fora do escopo declarado na Seção~\ref{sec:delimitacao}.

\citeonline{boehm1984} argumenta ainda que a verificação deve incidir
preferencialmente sobre os artefatos iniciais --- requisitos e projeto --- porque
são eles que determinam a maior parte do custo. Verificar apenas código é
verificar tarde.

\section{Testes e critérios de aceitação}
\label{sec:testes}

Diversas estratégias de teste coexistem. \citeonline{myers2011} enfatiza que o
objetivo do teste é \emph{encontrar defeitos}, não confirmar que o sistema
funciona, e que um caso de teste que não pode falhar não tem valor. Daí o
princípio de projetar casos a partir de condições limítrofes e de entradas
inválidas.

Duas práticas são diretamente relevantes para o processo proposto. A primeira é o
desenvolvimento orientado por testes, que inverte a ordem habitual: o teste que
expressa o comportamento desejado é escrito antes da implementação
(\citeonline{beck2002tdd}). A segunda é a especificação por exemplo, que
substitui descrições abstratas por exemplos concretos e executáveis do
comportamento esperado (\citeonline{adzic2011}). Ambas compartilham a mesma tese:
o critério de aceitação precisa existir antes da implementação, sob pena de se
tornar uma racionalização posterior do que foi produzido.

Neste trabalho, o critério de aceitação é o elo entre requisito e evidência. A
cobertura de verificação é definida como:

\begin{equation}
	A = \frac{\left| \{ \mathrm{CA}_i \mid \mathrm{resultado}(\mathrm{CA}_i) = \mathrm{PASS} \} \right|}
	         {\left| \{ \mathrm{CA}_i \} \right|}
	\label{eq:cobertura}
\end{equation}

\noindent em que $A$ é a fração de critérios de aceitação satisfeitos. A equação
mede conformidade com a especificação --- verificação, no sentido de
\citeonline{boehm1984} --- e não adequação ao uso.

\section{Qualidade de produto de software}
\label{sec:qualidade}

A ISO/IEC 25010 organiza a qualidade de produto em características como adequação
funcional, eficiência de desempenho, compatibilidade, usabilidade, confiabilidade,
segurança, manutenibilidade e portabilidade (\citeonline{iso25010}). O valor
prático desse modelo, para este trabalho, é oferecer vocabulário para classificar
requisitos não funcionais que, de outro modo, permaneceriam adjetivos soltos.

A Tabela~\ref{tab:mapeamento-iso} mapeia os requisitos não funcionais do alvo nas
características da norma. O mapeamento revela uma concentração: oito dos nove
requisitos não funcionais distribuem-se em apenas três características, o que
indica que a especificação priorizou fortemente portabilidade e adequação de
interface, em detrimento de eficiência de desempenho e confiabilidade.

\begin{table}[htb]
	\caption{Mapeamento dos requisitos não funcionais do alvo nas características de qualidade da ISO/IEC 25010}
	\label{tab:mapeamento-iso}
	\centering
	\footnotesize
	\begin{tabularx}{\textwidth}{@{}p{3.6cm} X@{}}
		\toprule
		\textbf{Característica (ISO/IEC 25010)} & \textbf{Requisitos não funcionais do alvo} \\
		\midrule
		Adequação funcional & NFR-008 (ausência de erros no carregamento) \\
		\addlinespace
		Eficiência de desempenho & --- (nenhum requisito não funcional dedicado) \\
		\addlinespace
		Compatibilidade & NFR-001 (sem dependências externas) \\
		\addlinespace
		Usabilidade & NFR-002 (layout centralizado), NFR-003 (tema escuro), NFR-007 (entrada não bloqueante), NFR-009 (movimento reduzido) \\
		\addlinespace
		Confiabilidade & --- (apenas tolerância a falha de armazenamento, tratada em FR-010) \\
		\addlinespace
		Segurança & --- (fora do escopo: aplicação local, sem rede nem dados pessoais) \\
		\addlinespace
		Manutenibilidade & NFR-005 (código comentado e organizado) \\
		\addlinespace
		Portabilidade & NFR-004 (responsividade), NFR-006 (nitidez em telas de alta densidade) \\
		\bottomrule
	\end{tabularx}
	\legend{Fonte: elaborado pelo autor (2026) a partir de \citeonline{iso25010} e da especificação do alvo.}
\end{table}

A ausência de requisito não funcional dedicado a eficiência de desempenho é uma
lacuna da especificação, não do produto: o alvo executa um laço contínuo de
renderização e, sem meta declarada de consumo ou de latência, não há critério
objetivo para reprová-lo nesse aspecto. Essa constatação reaparece na discussão
(Capítulo~\ref{cap:discussao}).

\section{Arquitetura e decisões de projeto}
\label{sec:arquitetura}

Arquitetura de software é o conjunto de decisões estruturais cujo custo de
mudança é alto (\citeonline{bass2021}). Documentar decisões arquiteturais
--- o problema enfrentado, a alternativa escolhida, as alternativas descartadas e
a justificativa --- preserva o raciocínio que deu origem ao projeto. Sem esse
registro, o código permanece como única fonte, e o código diz \emph{o que} foi
feito, não \emph{por que} nem \emph{em vez do quê}.

Padrões de projeto oferecem soluções nomeadas para problemas recorrentes
(\citeonline{gamma1994}), e a separação de responsabilidades é critério central
de qualidade estrutural. Neste trabalho, dois padrões são pertinentes: a
separação entre atualização de estado e apresentação, que isola a lógica de
simulação da lógica de desenho; e o uso de um objeto único de configuração, que
centraliza constantes e evita a duplicação de valores mágicos.

A integração contínua e a entrega contínua fornecem o contexto organizacional em
que portões de verificação automatizados se inserem
(\citeonline{humble2010}; \citeonline{forsgren2018}). O framework proposto no
Capítulo~\ref{cap:framework} é uma adaptação desses portões para um fluxo
executado por agentes.

\section{Modelos de linguagem de grande porte}
\label{sec:llm}

A arquitetura \textit{Transformer}, baseada em mecanismos de atenção, substituiu
as recorrências anteriores por processamento paralelo de dependências de longo
alcance (\citeonline{vaswani2017}). Sobre essa base, modelos com centenas de
bilhões de parâmetros demonstraram capacidade de executar tarefas a partir de
poucos exemplos fornecidos no próprio texto de entrada, sem re-treinamento
(\citeonline{brown2020}).

A capacidade que interessa a este trabalho não é a fluência, e sim o
\emph{raciocínio em etapas}. \citeonline{wei2022} mostraram que solicitar ao
modelo que explicite o raciocínio intermediário melhora o desempenho em tarefas
que exigem múltiplos passos. Isso importa porque verificação, análise de
requisito e decomposição de tarefa são, todas, atividades de múltiplos passos.

Duas limitações precisam ser declaradas, porque afetam a confiabilidade do
processo. A primeira é a geração de conteúdo plausível e incorreto, produzido com
o mesmo grau de confiança do conteúdo correto; é precisamente esse fenômeno que
torna indispensável exigir evidência independente da afirmação. A segunda é a
dependência do contexto: o comportamento do modelo é função do que está na janela
de entrada, de modo que informação ausente, obsoleta ou contraditória produz
saída degradada.

\section{Agentes de inteligência artificial}
\label{sec:agentes}

Um agente é um sistema que percebe um ambiente e age sobre ele
(\citeonline{russell2021}). A transposição dessa definição para modelos de
linguagem produziu agentes capazes de intercalar raciocínio e ação, alternando
etapas de deliberação com a execução de ferramentas e a observação do resultado
(\citeonline{yao2023react}). A capacidade de refletir sobre o próprio erro,
armazenando a avaliação crítica e reutilizando-a em tentativas seguintes,
melhora o desempenho em tarefas iterativas (\citeonline{shinn2023reflexion}).

Agentes introduzem dois problemas ausentes em modelos usados como meros
geradores de texto. O primeiro é a \emph{deriva de objetivo}: um agente que pode
escolher o próximo passo pode escolher um passo que não serve ao objetivo
original. O segundo é a \emph{opacidade}: a sequência de decisões que levou ao
artefato final pode não ser reconstruível. Ambos são problemas de controle, e é
para eles que o framework do Capítulo~\ref{cap:framework} se orienta.

\section{Sistemas multiagente}
\label{sec:multiagente}

Sistemas multiagente clássicos estudam coordenação, cooperação e negociação entre
agentes autônomos (\citeonline{wooldridge2009}). Aplicada a modelos de linguagem,
a ideia deu origem a frameworks que distribuem papéis especializados. MetaGPT
organiza papéis análogos a uma equipe de software e representa os artefatos
intermediários de forma estruturada (\citeonline{hong2024metagpt}). ChatDev
modela o desenvolvimento como diálogo entre agentes em papéis de projetista,
programador e revisor (\citeonline{qian2024chatdev}). AutoGen generaliza a
conversação entre agentes configuráveis como primitiva de construção
(\citeonline{wu2023autogen}).

Esses trabalhos tratam principalmente da \emph{geração cooperativa}. A
contribuição pretendida neste trabalho é complementar: trata-se de fixar os
\emph{artefatos de controle} --- especificação, decisão, tarefa, evidência --- e
as condições sob as quais o fluxo deve parar por falta de informação. A
comparação detalhada está no Capítulo~\ref{cap:relacionados}.

\section{Contexto, custo e disciplina documental}
\label{sec:contexto-custo}

Como o comportamento do agente é função da janela de contexto, o volume de
documentação carregada tem custo duplo: custo computacional e custo de ruído.
Informação irrelevante não é neutra --- ela compete por atenção com a informação
relevante. Daí decorre um princípio de projeto adotado neste trabalho:
documentos de controle devem ser curtos, e a leitura do repositório deve ser
incremental e orientada à tarefa corrente, não exaustiva.

Em tensão com esse princípio está a exigência de rastreabilidade, que empurra na
direção oposta, a de mais documentação. O framework proposto resolve a tensão por
especialização: cada informação tem \emph{um} documento dono, e os demais
apontam para ele em vez de repeti-lo. Uma matriz de rastreabilidade existe para
conectar, não para duplicar.

\section{O domínio de aplicação}
\label{sec:dominio}

O estudo de caso é um jogo digital de movimentação em grade, entregue como página
web única. Três características do domínio tornam-no adequado ao propósito.

\textbf{Estado discreto e verificável.} O jogo é uma máquina de estados finitos
sobre uma grade discreta, com condições de término bem definidas (colisão com
borda ou com o próprio corpo). Isso permite afirmar propriedades precisas e
verificá-las por observação, sem depender de julgamento subjetivo.

\textbf{Lógica de tempo real.} O jogo exige um laço contínuo de atualização e
desenho. A plataforma web oferece \texttt{requestAnimationFrame} para
sincronizar o desenho com a taxa de atualização do dispositivo
(\citeonline{mdn2024raf}), e o desacoplamento entre taxa de simulação e taxa de
quadros é uma decisão de projeto com consequências mensuráveis
(Capítulo~\ref{cap:implementacao}).

\textbf{Apresentação e persistência locais.} A superfície de desenho vetorial
bidimensional é adequada a grade discreta (\citeonline{mdn2024canvas}), e o
armazenamento local do navegador permite persistir estado entre sessões sem
servidor (\citeonline{mdn2024storage}). A acessibilidade é orientada pelas
diretrizes WCAG (\citeonline{w3c2023wcag22}), e a linguagem é regida pelo padrão
ECMAScript (\citeonline{ecma2023}) sobre a especificação viva de HTML
(\citeonline{whatwg2024html}).

\section{Síntese e lacuna identificada}
\label{sec:sintese}

A fundamentação apresentada sustenta três conclusões que orientam o restante do
trabalho.

Primeira: a verificação depende de critérios definidos \emph{antes} da
implementação; critérios escritos depois são descrição, não verificação.

Segunda: requisitos não funcionais só são verificáveis quando convertidos em
métricas; quando permanecem adjetivos, eles não entram em nenhum portão de
aceitação, e a lacuna passa despercebida justamente porque não há teste que possa
falhar.

Terceira: agentes de linguagem são úteis para gerar e para revisar, mas não são
fonte de evidência. A evidência precisa vir de execução observável e
reproduzível.

A lacuna que o trabalho aborda é a ausência, na literatura de agentes LLM para
engenharia de software, de um processo que trate explicitamente os artefatos de
controle e as condições de parada, com ênfase em verificabilidade medida. O
Capítulo~\ref{cap:relacionados} sustenta essa afirmação por comparação.

A verificação, por sua vez, não é um evento único e sim um conjunto de camadas
aplicadas em momentos distintos do trabalho. A Figura~\ref{fig:verificacao}
representa essas camadas e registra quais foram efetivamente aplicadas na
produção deste documento. A distinção entre camadas aplicadas e camadas não
executadas é parte do resultado, não uma omissão.

%% ---------------------------------------------------------------------
%%  FIGURA - Camadas de revisao da monografia (DG-10)
%% ---------------------------------------------------------------------
\begin{figure}[htb]
	\caption{Camadas de revisão aplicadas à produção da monografia}
	\label{fig:verificacao}
	\centering
	\begin{tikzpicture}[
			font=\footnotesize,
			cam/.style={draw, rounded corners=2pt, minimum width=8.2cm, minimum height=0.7cm, align=center, fill=black!4},
			>=Stealth,
			fl/.style={->, thick}
		]

		\node[cam] (r1) at (0,0) {1. Conteúdo};
		\node[cam, below=0.2cm of r1] (r2) {2. Engenharia};
		\node[cam, below=0.2cm of r2] (r3) {3. Evidências};
		\node[cam, below=0.2cm of r3] (r4) {4. Coerência};
		\node[cam, below=0.2cm of r4] (r5) {5. Bibliográfica};
		\node[cam, below=0.2cm of r5] (r6) {6. ABNT};
		\node[cam, below=0.2cm of r6] (r7) {7. Linguística};
		\node[cam, below=0.2cm of r7] (r8) {8. LaTeX};
		\node[cam, below=0.2cm of r8] (r9) {9. Inspeção do PDF};

		\draw[fl] (5.2,0.0) -- (5.2,-7.4);
	\end{tikzpicture}
	\legend{Fonte: elaborado pelo autor (2026). Neste trabalho, as camadas 1 a 5 foram aplicadas durante a
	produção; a camada 6 foi parcialmente aplicada e depende do manual institucional; as camadas 8 e 9
	\textbf{não foram executadas} por ausência de compilador LaTeX no ambiente --- ver
	\texttt{MONOGRAFIA\_STATUS.md}, seção~3.}
\end{figure}

